\section{From spherical arrangements to bounded path systems}
\label{sec:planar}

Theorem~\ref{thm:bounded-patterns} excludes path systems whose number of
paths, total length in units of $L$, and number of comparison intervals
are bounded.  The topological argument will produce a planar pairing of
paths with none of those bounds.  We show that every such pairing
yields a system covered by the theorem.  Planarity has two roles: an
Euler characteristic count organizes the pairing into few intervals,
and a separator reduces the number of paths while losing few paired
positions.  This is the spherical-arrangement method of
\cite[Section ``From bounded patterns to spherical arrangements'']{OpenAI2026ZD};
we include the geometric and quantitative details needed here.

\begin{definition}[Reduced spherical arrangement]
\label{def:reduced-arrangement}
A \emph{reduced spherical arrangement over $\Gamma$} consists of
finitely many pairwise disjoint piecewise linear closed disks in an
oriented sphere, together with the following data.
Each boundary reads a positive-length cyclically immersed closed path
in $\Gamma$, except possibly one boundary.  That boundary reads a
positive-length immersed interval path starting at $x_A$ or $x_B$,
with its two endpoints represented at a single marked boundary point,
called the \emph{break}.  No immersion condition is imposed across the
break.  At least one ordinary, cyclically immersed boundary is required.

The paths are read in the boundary orientations induced by the disks.
Their edge occurrences are paired in distinct twos by mutually disjoint
piecewise linear arcs.  The interiors of the arcs lie outside the disks,
and each occurrence contains exactly one arc endpoint in its interior.
Paired occurrences have inverse signed letters and use distinct
underlying undirected edges of $\Gamma$.
\end{definition}

The distinction between occurrences and graph edges is essential:
different occurrences may traverse the same edge, but the two ends of
any one pairing arc may not.  The graph obtained by collapsing each
disk to a vertex may be disconnected and may have loops or parallel
edges.

Proposition~\ref{prop:cone-criterion} will show that the absence of
these arrangements makes the cone complex aspherical and protects the
two roots. The present section proves that the random graphs can be
chosen to have this absence property.

\begin{theorem}[Absence of reduced spherical arrangements]
\label{thm:no-arrangements}
In the random matching model of Section~\ref{sec:random}, conditioned
on girth at least $L$, the probability that $\Gamma$ admits a reduced
spherical arrangement tends to zero as $m\equiv1\pmod4$ tends to
infinity.  Consequently, for every sufficiently large admissible $m$
there is a matching outcome with girth at least $L$, the component
diameter bound of Proposition~\ref{prop:diameter}, and no reduced
spherical arrangement.
\end{theorem}

We first establish the three deterministic tools used in the proof.
The first closes path segments without changing any original
occurrences.  The other two depend only on planarity.

\subsection{Closing immersed segments}

Fix the constant $d$ of Proposition~\ref{prop:diameter}, and work for
now with an outcome in which every component of $\Gamma$ has diameter
at most $dL$.  Each side of $\Gamma$ has at most $n$ vertices and minimum
degree at least $129$.

\begin{lemma}[Short closures]
\label{lem:closure}
There is a fixed integer $D_0$, depending only on $d$ and $c_0$, with
the following property for all sufficiently large admissible $n$.
If every component of $\Gamma$ has diameter at most $dL$, then every
positive-length immersed path segment in $\Gamma$ extends to a
cyclically immersed closed path by appending at most $D_0L$ edges.
The original segment is retained without cancellation.  One may take
any integer
\[
 D_0\ge d+\frac{8}{c_0\log2}+8.
\]
\end{lemma}

\begin{proof}
We first construct a short based loop that avoids specified edges at
its basepoint.  For large $n$ we have $L\ge3$, so the girth condition
excludes loops and parallel edges.  Delete at most two prescribed
undirected edges incident to a vertex $z$.  In the remaining graph
every vertex still has degree at least $127$, and in particular at
least three.  With
\[
 k=\lceil\log_2n\rceil+1,
\]
there are at least $3\cdot2^{k-1}>n$ non-backtracking walks of length
$k$ starting at $z$.  Two distinct such walks, say $P$ and $Q$, end
at the same vertex.  Reduce $PQ^{-1}$ as a linear path.  Cancellation
removes their common terminal segment.  Since $P$ and $Q$ are distinct
and have equal length, it does not remove either entire walk.  We
obtain a nonempty linearly immersed based loop of length at most $2k$;
its first and last edges avoid the deleted edges.  The loop need not
be cyclically immersed at $z$.

Let the given segment $S$ run from $a$ to $b$, and choose a geodesic
$R$ from $b$ to $a$, of length at most $dL$.  If $R$ is nonempty,
construct a loop $J_b$ at $b$ avoiding the last edge of $S$ and the
first edge of $R$, and a loop $J_a$ at $a$ avoiding the last edge of
$R$ and the first edge of $S$.  Then
\[
 S J_b R J_a
\]
is cyclically immersed: each constituent is immersed as a linear path,
and each of the four joins avoids immediate reversal.  If $R$ is empty,
then $a=b$, and one loop at $a$ avoiding the first and last edges of
$S$ makes $SJ_a$ cyclically immersed.  Thus the total added length is
at most $dL+4k$ in either case.  Finally,
$L\ge(c_0/2)\log n$ and $L\ge1$ for large $n$, so
\[
 k\le\left(\frac{2}{c_0\log2}+2\right)L.
\]
This gives the stated choice of $D_0$.
\end{proof}

\subsection{Organizing the pairing into intervals}

We next count how often a planar pairing can change its consecutive
partner interval.  A pair of neighboring arc ends usually has another
neighboring pair opposite it across a digon.  Only the other gaps,
and the places where we choose to split a path, require interval cuts.
This is the linear comparison-interval argument of
\cite[Lemma~4.3]{OpenAI2026ZD}.

\begin{lemma}[Paired intervals]
\label{lem:paired-intervals}
Suppose a reduced spherical arrangement with $N_0$ disk boundaries is
divided into $N'\ge N_0$ pieces by cutting its paths into consecutive
segments.  An intact path counts as one piece, and no piece crosses
the exceptional break.  Then all original occurrence pairings can be
represented by at most $6N'$ disjoint pairs of intervals.  Each interval
is contained in one designated piece, and each pair compares inverse
strings in reverse boundary order.  The intervals can also be required
not to cross one prescribed indexing gap between occurrences on each
original path, using the break as the exceptional indexing gap.
\end{lemma}

\begin{proof}
Thicken the disks and the pairing arcs to a closed regular neighborhood
$R$ made of disks and rectangular bands.  Write $E$ for the number of
bands, or equivalently the number of pairing arcs.  A \emph{gap} is the
portion of an original disk boundary between successive arc ends in
cyclic order.  After thickening, each gap gives a corner of a
complementary region.  For each connected component $Q$ of
$S^2\setminus\operatorname{int}R$, let $\ell_Q$ be the number of band
sides along its boundary, counting incidences.  This is also its number
of gap incidences.  Every disk has positive boundary length and every
occurrence is paired, so every boundary component of $R$ encounters a
band and $\ell_Q\ge1$.

The neighborhood $R$ retracts to the graph with $N_0$ vertices and $E$
edges.  Additivity of Euler characteristic when surfaces are glued
along boundary circles therefore gives
\[
 \chi(R)=N_0-E,\qquad
 \sum_Q\chi(Q)=2-N_0+E,\qquad
 \sum_Q\ell_Q=2E.
\]
In particular,
\begin{equation}
 \sum_Q\bigl(\ell_Q-2\chi(Q)\bigr)=2N_0-4.
 \label{eq:arrangement-euler}
\end{equation}
This identity does not require $R$ to be connected; the regions $Q$
may have several boundary components.

Call a gap \emph{good} if its region $Q$ is a disk digon whose two sides
belong to distinct bands.  The two corners of such a digon are opposite
good gaps.  Figure~\ref{fig:good-digon} shows how coherent boundary
orientations reverse the order of occurrences across such a digon.

\begin{figure}[tbp]
\centering
\begin{tikzpicture}[x=1cm,y=1cm,font=\small,
  line cap=round,line join=round,>=Latex]
  \fill[black!10] (0.60,2.20) rectangle (6.35,2.83);
  \fill[black!10] (0.60,-0.23) rectangle (6.35,0.40);
  \fill[black!20] (2.03,0.40) rectangle (2.37,2.20);
  \fill[black!20] (4.63,0.40) rectangle (4.97,2.20);
  \draw (0.60,2.20) -- (6.35,2.20);
  \draw (0.60,0.40) -- (6.35,0.40);
  \draw (2.03,0.40) -- (2.03,2.20);
  \draw (2.37,0.40) -- (2.37,2.20);
  \draw (4.63,0.40) -- (4.63,2.20);
  \draw (4.97,0.40) -- (4.97,2.20);
  \draw[->,line width=0.8pt] (5.47,2.20) -- (6.07,2.20);
  \draw[->,line width=0.8pt] (6.07,0.40) -- (5.47,0.40);
  \node at (2.20,2.55) {$e$};
  \node at (4.80,2.55) {$f$};
  \node at (2.20,0.04) {$e'$};
  \node at (4.80,0.04) {$f'$};
  \node at (3.50,1.30) {$Q$};
  \node[font=\footnotesize] at (3.50,0.91) {disk digon};
  \draw[line width=1pt] (3.50,2.06) -- (3.50,2.34);
  \draw[line width=1pt] (3.50,0.26) -- (3.50,0.54);
  \node[font=\footnotesize] at (3.50,2.57) {good gap};
  \node[font=\footnotesize] at (3.50,0.02) {opposite gap};
  \node[font=\footnotesize,rotate=90] at (2.20,1.30) {band};
  \node[font=\footnotesize,rotate=90] at (4.80,1.30) {band};
  \node[align=left,anchor=west] at (7.02,1.90)
    {Coherent boundary orders:\\$e,f$\quad and\quad $f',e'$};
  \node[align=left,anchor=west,font=\footnotesize] at (7.02,0.74)
    {The ticks mark opposite gaps.\\A cut at one is copied to the other.};
\end{tikzpicture}
\caption{A good gap and its opposite gap, shown locally between two
distinct bands.  The complementary region $Q$ is a disk digon.  The
coherent boundary orientations reverse the order of the paired
occurrences: $e,f$ on one boundary correspond to $f',e'$ on the other.
Thus an indexing or splitting cut at a good gap must also cut its
opposite gap, so that the remaining runs pair as whole intervals.
The symbols denote occurrences, not underlying edges of $\Gamma$;
shading only distinguishes disk portions and bands in the surface.}
\label{fig:good-digon}
\end{figure}


Let $B$ be the number of gaps that are not good.  We bound $B$ by
treating the possible complementary regions as follows.

A disk monogon has one gap between two occurrences joined by the same
pairing arc.  Those neighboring letters are inverse.  Immersion rules
this out except at the exceptional break, so the number $M$ of disk
monogons is at most one.  Their total contribution to
Equation~\ref{eq:arrangement-euler} is $-M$.

A disk digon using the same band twice turns immediately back along
that band at an endpoint.  The associated gap is the sole gap at a
disk with only one incident arc end: otherwise another band end would
intervene in its cyclic order.  Different such digons use different
degree-one gaps.  Hence there are at most $N_0$ such digons and at
most $2N_0$ gaps on them.

Every remaining region is either a disk with $\ell_Q\ge3$, or a
connected planar surface with at least two boundary components.  In
the first case $\ell_Q\le3(\ell_Q-2)$; in the second case
$\chi(Q)\le0$, and the stronger inequality
$\ell_Q\le\ell_Q-2\chi(Q)$ holds.  These regions have positive
contributions to Equation~\ref{eq:arrangement-euler}, with total
$2N_0-4+M$.  Digons contribute zero.  It follows that
\begin{equation}
 B\le M+2N_0+3(2N_0-4+M)\le8N_0.
 \label{eq:non-good-gaps}
\end{equation}

Across a good digon, the partners of consecutive occurrences $e,f$
occur consecutively in the opposite order $f',e'$.  Taking opposite
corners is thus an involution on good gaps.  Cut every original
boundary at every non-good gap.  Also mark its prescribed indexing
gap, using the break on the exceptional boundary, and mark all gaps
used in dividing the paths into the designated pieces.  These extra
marks number at most $N_0+N'$.  For each marked good gap, cut both it
and its opposite gap.  The resulting cut set is stable under the
good-gap involution, and its cardinality $J$ satisfies
\[
 J\le B+2(N_0+N')\le12N'.
\]

Every boundary has a cut, so the remaining consecutive runs of
occurrences are linear intervals.  All of them lie within designated
pieces.  Pairing reverses each uncut adjacency, and stability of the
cuts ensures that it takes a whole maximal run onto another whole
maximal run.  A run cannot be paired with itself.  For an odd-length
run this would fix its middle occurrence, whereas occurrences are
paired in distinct twos.  For an even-length run it would pair the
two middle occurrences and make their adjacent letters inverse.
That contradicts immersion inside the run; the exceptional break has
been cut.  The $J$ runs therefore pair in distinct twos, giving at most
$J/2\le6N'$ interval pairs as claimed.
\end{proof}

\subsection{A recursive separator bound}

Once a long boundary has been divided into pieces of controlled
length, deleting a small fraction of the pieces loses only a small
fraction of all occurrences.  The following consequence of the planar
separator theorem, also used in \cite[Lemma~4.4]{OpenAI2026ZD}, makes
the remaining groups of pieces bounded in size.

\begin{lemma}[Recursive planar separation]
\label{lem:separator}
Let
\[
 C_{\mathrm{sep}}=\frac{2\sqrt2}{1-\sqrt{2/3}}.
\]
For every integer $K_0\ge1$, a finite planar multigraph with $N'$
vertices has a set of at most $C_{\mathrm{sep}}N'/\sqrt{K_0}$ vertices
whose deletion leaves every connected component with at most $K_0$
vertices.
\end{lemma}

\begin{proof}
Deleting loops and replacing parallel edges by single edges preserves
the components after every possible vertex deletion, so it suffices
to consider a simple planar graph.  The Lipton--Tarjan separator theorem
\cite{LiptonTarjan1979} gives, for a simple planar graph on $s$ vertices,
a set of at most $2\sqrt{2s}$ vertices whose removal leaves components
with at most $2s/3$ vertices.  Apply this theorem recursively to each
component with more than $K_0$ vertices, starting separately on the
original connected components.

At a recursive call on $s>K_0$ vertices, charge each of those vertices
$2\sqrt2/\sqrt{s}$.  The total charge at that call is $2\sqrt{2s}$ and
bounds the number of vertices removed there.  Follow any fixed original
vertex through the calls containing it, until it is deleted or belongs
to a component of size at most $K_0$.  The successive charged component
sizes decrease by a factor at most $2/3$.  Reading backwards from the
last charged size, which is greater than $K_0$, bounds its total charge
by
\[
 \frac{2\sqrt2}{\sqrt{K_0}}
 \sum_{j\ge0}(2/3)^{j/2}
 =\frac{C_{\mathrm{sep}}}{\sqrt{K_0}}.
\]
This applies also to a vertex deleted by a separator.  Summing the
charges over all $N'$ original vertices bounds the total number
deleted and proves the lemma.
\end{proof}

\subsection{Extracting one bounded path system}

We now combine these tools.  The original paths will first be divided
and closed, then separated into clusters.  We retain a cluster whose
proportion of unpaired occurrences and number of interval pairs are
both small.  Controlling the total losses, rather than those of every
cluster, is enough to find one such cluster.

\begin{proof}[Proof of Theorem~\ref{thm:no-arrangements}]
Let $\varepsilon>0$ be the constant in
Theorem~\ref{thm:bounded-patterns}, independent of its fixed bounds
$K_0,C,I$.  Fix $d$ from Proposition~\ref{prop:diameter} and $D_0$ from
Lemma~\ref{lem:closure}. We will cut long boundaries into pieces with
at most $UL$ original occurrences, pay at most $D_0L$ to close each
cut piece, and delete at most an $\eta$ fraction of the pieces.
We choose $U$ large enough to make the closure overhead small and
$\eta$ small enough to control the occurrences lost in separation.
Specifically, choose the following constants in order:
\begin{align}
 U&\ge\max\{3,\lceil96D_0/\varepsilon\rceil\},
       &&\text{$U$ an integer},\notag\\
 \eta&=\min\{1/(32U),\varepsilon/(64U)\},\qquad
 K_0=\max\{1,\lceil(C_{\mathrm{sep}}/\eta)^2\rceil\},
 \label{eq:extraction-constants}\\
 C&=K_0(U+D_0),\qquad I=\lceil192C\rceil.\notag
\end{align}
In particular,
\begin{equation}
 \frac{3D_0}{U}\le\frac{\varepsilon}{32},\qquad
 2\eta U\le\frac1{16},\qquad
 2\eta U\le\frac{\varepsilon}{32}.
 \label{eq:extraction-losses}
\end{equation}
All these constants are fixed before $n=vm$ tends to infinity through
$m\equiv1\pmod4$.  Their order of dependence is
\[
 (\text{types},c_0,\delta)\ ;\quad
 (\varepsilon,d)\ ;\quad D_0\ ;\quad U\ ;\quad\eta\ ;\quad
 K_0\ ;\quad (C,I)\ ;\quad n.
\]

Suppose there is a reduced spherical arrangement on an outcome
satisfying the diameter bound.  Let $N_0$ be its number of boundaries
and $H_0$ their total length.  An ordinary boundary is present, and
girth gives $H_0\ge L$.

Keep every path of length at most $UL$ intact.  For a longer path of
length $h$, put $r=\lceil h/(UL)\rceil$ and divide it into $r$
consecutive segments whose integer lengths differ by at most one.
For large $L$ each segment has length between $UL/3$ and $UL$:
$h/r>UL/2$, so $\lfloor h/r\rfloor\ge UL/3$ when $UL\ge6$, while
$\lceil h/r\rceil\le UL$.  Use the linear order from the root on an
exceptional path, so no segment crosses its break.  Close each segment
of every path that was split using Lemma~\ref{lem:closure}.  Every
newly appended occurrence is left unpaired; intact paths receive no
added edges.

Let $N'$ be the number of resulting paths, and let $H_{\mathrm{add}}$
be their total added length.  Each resulting path contains at most
$UL$ original occurrences and has total length at most $(U+D_0)L$.
Since a split path contributes at most $3h/(UL)$ pieces,
\begin{equation}
 H_{\mathrm{add}}\le\frac{3D_0}{U}H_0.
 \label{eq:added-length}
\end{equation}
Every piece has at least $L$ original occurrences, by girth for intact
ordinary paths and by $UL/3\ge L$ for split paths, except possibly the
one intact exceptional path.  Thus
\begin{equation}
 N_0\le N',\qquad N'L\le H_0+L\le2H_0.
 \label{eq:extraction-piece-count}
\end{equation}
If the exceptional path was split, all its pieces have become closed
cyclically immersed paths.  Otherwise it remains a single immersed
interval starting at its original root.

Apply Lemma~\ref{lem:paired-intervals} before the closures are appended.
The original pairing is specified by at most $6N'$ interval pairs,
each lying in its designated pieces.  Choose the indexing points of
intact ordinary paths as in that lemma; index every newly closed path
with its original segment followed by its closure.  The interval
comparisons therefore remain valid in the indexed resulting paths.

Form a multigraph whose vertices are the $N'$ pieces and whose edges
are the original pairing arcs.  It is planar: within each old disk,
replace its vertex by one vertex for each consecutive block of arc
ends belonging to a piece.  Each block can be connected to its new
vertex inside a separate sector of the disk.  This produces disjoint
connections, also for original arcs with both ends at the same disk.
The appended closures have no pairing arcs and play no role in this
embedding.

By Lemma~\ref{lem:separator} and the choice of $K_0$, delete at most
$\eta N'$ vertices, leaving components with at most $K_0$ pieces.
Call these components \emph{clusters}.  If $M_{\mathrm{del}}$ is the
number of original occurrences in deleted pieces, then
\begin{equation}
 M_{\mathrm{del}}\le\eta N'UL\le2\eta U H_0.
 \label{eq:deleted-occurrences}
\end{equation}
Keep an original pairing precisely when both of its pieces survive.
At most $M_{\mathrm{del}}$ surviving original occurrences lose their
partners, because each deleted occurrence had just one partner.
An interval pair is either wholly retained or wholly discarded:
each of its two intervals lies in one piece.  Every retained pair
belongs to a single cluster.

For each cluster $Z$, denote its total path length by $H_Z$, including
added closures, its number of unpaired occurrences by $b_Z$, and its
number of retained interval pairs by $I_Z$.  Equations~\ref{eq:added-length}--\ref{eq:deleted-occurrences}
and the parameter choices give
\begin{align}
 \sum_Z H_Z&\ge H_0-M_{\mathrm{del}}\ge15H_0/16,
 & H_Z&\le K_0(U+D_0)L=CL,\notag\\
 \sum_Z b_Z&\le M_{\mathrm{del}}+H_{\mathrm{add}}
                \le\varepsilon H_0/16,
 &\sum_Z I_Z&\le6N'\le12H_0/L.
 \label{eq:cluster-budgets}
\end{align}
These are totals over all remaining clusters; an individual cluster
need not yet satisfy the bounded-pattern conditions.

Discard all clusters with $b_Z>\varepsilon H_Z$.  Their total length
is at most $\varepsilon^{-1}\sum_Z b_Z\le H_0/16$.  Also discard all
clusters with $I_Z>192H_Z/L$.  Their total length is at most
$(L/192)\sum_Z I_Z\le H_0/16$.  After these two exclusions, at least
$13H_0/16$ of cluster length remains.

A cluster with $H_Z<L$ contains no positive-length cyclically immersed
path, by girth.  It can therefore consist only of the one intact
exceptional path.  There is at most one such cluster.  If its length
is $h<L$, the original arrangement also has an ordinary boundary of
length at least $L$, so $H_0\ge h+L>2h$.  Discarding this final cluster,
if present, costs less than $H_0/2$.  A positive total length remains,
since $13H_0/16-H_0/2=5H_0/16>0$.

Choose one remaining cluster.  It has at most $K_0$ positive-length
paths and satisfies
\[
 L\le H_Z\le CL,\qquad b_Z\le\varepsilon H_Z,\qquad
 I_Z\le192H_Z/L\le192C\le I.
\]
All paths are cyclically immersed except possibly the entire intact
exceptional interval, still starting at $x_A$ or $x_B$.  The retained
comparisons still pair distinct underlying edges.  Closures may
revisit edges or meet other paths, but their occurrences are all
unpaired, which Theorem~\ref{thm:bounded-patterns} permits.  This
cluster is therefore a path system excluded by that theorem for the
single fixed triple $(K_0,C,I)$ in
Equation~\ref{eq:extraction-constants}.

The construction applies to every finite arrangement, with no bound on
$N_0$ or $H_0$, and always yields a system for that same triple.
The probability of such a system is $o(1)$ by
Theorem~\ref{thm:bounded-patterns}; the probability that the diameter
bound fails is $o(1)$ by Proposition~\ref{prop:diameter}.  This proves
the asserted exclusion probability.  The simultaneous girth,
diameter, and exclusion requirements consequently have positive
probability for all sufficiently large admissible $m$.
\end{proof}
